\scanpage{631}
\mspage{19}
\deleted{factorisons} posons

\[
(\qquad f=u_{\deleted{\unsure{!|*}}}^{*}(f'),
\]

et factorisons $f$ en $pi$, avec $i\in TC$, $p\in F$,
d'où

\[
u_!u^{\unsure{*|u}}(f')=u_{\deleted{\unsure{*}}!}(f)=u_!(p)\,u_!(i).
\]

Mais $u^*$ étant \unsure{pl.} fidèle $u_!u^*\xrightarrow{\sim}\mathrm{id}$, donc
on trouve une factorisation de \deleted{$\unsure{u_!(f)}$}
$f'$ en $u_!(p)\,u_!(i)$. \deleted{Comme $u_!(F)\subset F'$} Il suffit
donc de voir que

\[
u_!(p)\in F',\qquad u_!(i)\in W',
\]

ce qui équivaut à

\[
u^*u_!(p)\in F,\qquad u^*u_!(i)\in W
\]

et résulte des hypothèses \deleted{\unsure{b}} que $u^*u_!(F)\subset F$
et $u^*u_!(TC)\subset TC\subset W$, qed.

\uline{NB} J'ai noté les cas A) et C) par acquit de
conscience, sans application \unsure{précise|pratique} en vue. Le cas
\inserted{le plus} utile sera \deleted{\unsure{le}}\inserted{\unsure{justiciable du}} cas B), \deleted{qui}\inserted{et} mérite d'être
\unsure{réconforté|récapitulé} en un

% slanted
\margin{cf Scholie 2.10, qui devrait précéder le th. 2.3.}
\uline{Théorème} 2.3. Soit $(W,C,F)$ un triple de Quillen
spécial \unsure{dans} $M$. Pour \unsure{tt} $I$ dans $\mathrm{Cat}$, considérons
le triple $(W(I),C\langle I\rangle,F(I))$\inserted{\deleted{\unsure{dans}} $M(I)=\underline{\mathrm{Hom}}(I^{\mathrm o},M)$}, où les notations
sont celles des pages 10, donc $W(I),F(I)$ et
$TF(I)=\deleted{\unsure{T}}F(I)\cap W(I)$ se définissant fibre par fibre,

\[
C\langle I\rangle=(F(I))_{\deleted{\unsure{*}}}^{*}=\big(\unsure{v|u}_{I!}(C(I_0))\big)^{\sim}.
\]
% slanted
\margin{où $v_I:I_0\to I$ ($I_0$ cat. discrète \inserted{ayant \unsure{$=$|mêmes} objets} \deleted{\unsure{de cat.} \illegible{1}} \unsure{que} $I$) \deleted{\unsure{source de $I$}} et}
\deleted{Alors}

a) Pour \unsure{tte} $I$ dans $\mathrm{Cat}$, $(W(I),C\langle I\rangle,F(I))$ est
% archive sheet number: 630
